% !TEX root = main.tex

\section{Introduction}

\noindent Measuring the complexity of real numbers is of major importance in computer science. Consider a non-computable real number $s$, i.e., a real number which cannot be stored on a computer \cite[Section 10]{copeland2004computable}. One can store only an approximation of $s$, for instance by considering a finite length bitstring $y$ representing a prefix $x$ of the binary expansion $\xf$ of $s$. For a fixed approximation error $\varepsilon>0$, the required length of $y$, as a function of $\varepsilon$, depends on the \textit{algorithmic complexity} of the prefix $x$ of the binary expansion achieving error $\varepsilon$. The \textit{algorithmic complexity} of a binary sequence $x$, often referred to as \textit{Kolmogorov complexity}, is the length of the smallest binary sequence $y$, for which there exists an algorithm, such that when presented with $y$ as input, delivers $x$ as output \cite[Definition 10.1]{balcazar2012structural}\cite[Section 14.2]{cover1999elements}\cite[Definition 2.1.2]{li2008introduction}. The algorithmic complexity of the binary expansions of real numbers has been widely studied, but the algorithmic complexity of expansions in bases other than $2$ remains poorly understood. Several papers have established an equivalence between the algorithmic complexity of the expansions in different bases $q \in \N$ \cite[Theorem 6.1]{calude2005randomness}\cite[Theorem 5.1]{hertling1998randomness}\cite[Theorem 3]{staiger2002kolmogorov}. Here, we study the algorithmic complexity of expansions in noninteger bases, which display a much more sophisticated behavior. This type of expansions are often referred to as $\beta$-expansions \cite[(1)]{parry1960beta}\cite[Section 4]{renyi1957representations}. $\beta$-expansions display some redundancy properties that are used to design robust A/D converters \cite{daubechies2002beta}\cite{daubechies2006imperfect}\cite{daubechies2010GoldenRatioEncoder}. However, we discover that the algorithmic complexity of $\beta$-expansions can be much larger that the algorithmic complexity of binary expansions of the same number. This engenders problems of compressibility of the sequences generated by the aforementioned A/D converters. Fortunately, we find a fast algorithm to fix this issue, converting any $\beta$-expansion of potentially large algorithmic complexity into another $\beta$-expansion which algorithmic complexity is equal to the algorithmic complexity of the binary expansions.

% , has been studied widely in the literature in the context of dynamical systems \cite[Example 1.2.3]{dajani2002ergodic}\cite{dajani2002greedy}\cite{dajani2003random}. 

% \noindent Measuring the complexity of real numbers is of major importance in computer science, for the purpose of knowing which computations are allowed. Consider a non-computable  real number $s$, i.e. a real number which cannot be stored on a computer. We can store only an approximation of $x$, for instance by considering a finite bitstring representing a finite prefix of its binary expansion. For a fixed approximation error $\epsilon>0$, the size of this finite bitstring is dependent on the algorithmic complexity of the finite prefixes of the binary expansion of $s$. The algorithmic complexity of a binary sequence $x$, often referred to as Kolmogorov complexity, is the length of the smallest binary sequence $x'$, for which there exists an algorithm, such that when presented with $x'$ as input, it outputs $x$. The algorithmic complexity of the binary expansion of real numbers is widely studied, but the algorithmic complexity of other ways of representing real numbers remains poorly reported. However, knowing the algorithmic complexity of different representations may allow to define new and more efficient strategies to represent real numbers. To our knowledge, \cite{staiger2002kolmogorov} is the only paper to make a step in this direction: it establishes an equivalence between the algorithmic complexity of the $q$-ary expansions, with $q \in \N$, $q \geq 2$, i.e. representations of real numbers in any integer base. In this paper, we study the algorithmic complexity of the so-called $\beta$-expansions, which are representations of real numbers in a base $\beta \in (1,2)$. $\beta$-expansions display very interesting properties, one of which being redundancy: a given real number $s$ and a fixed $\beta \in (1,2)$, there can be infinitely many different $\beta$-expansions of $s$. In \cite{daubechies2002beta}, this redundancy is used to design robust A/D converters. However, nothing is known about the algorithmic complexity of the generated sequences, and they may potentially be way more complex than the binary expansion, displaying a robutness-complexity trade-off. In this paper, we show that for a given real number $s$, the binary expansion is a minimizer of algorithmic complexity, and that for every given $\beta \in (1,2)$, there exists a $\beta$-expansion of $s$ which achieves the lower bound of algorithmic complexity displayed by the binary expansion of $s$.